\begin{abstract}
War et al. (2025), ``Detection of Security Smells in IaC Scripts through Semantics-Aware Code and Language Processing'' (University of Luxembourg, arXiv:2509.18790; a preprint whose peer-review status could not be confirmed), fine-tune CodeBERT/LongFormer to detect misconfigurations in Ansible/Puppet scripts, substantially beating prior TF-IDF-based baselines. Their own ablation study is the source of the gap this project targets: removing natural-language context (comments, task descriptions) causes precision to collapse sharply while recall holds or rises -- CodeBERT on Ansible drops from precision 0.92 to 0.46 (F1 0.90 to 0.58). This project reproduces that qualitative collapse on a synthetic IaC benchmark deliberately split into "easy" misconfigurations (safe/unsafe distinguishable from code tokens alone) and "hard" ones (the code is lexically identical either way; only an explanatory comment reveals the label) -- a lightweight TF-IDF + logistic regression classifier drops from precision/recall/F1 of 1.00/1.00/1.00 with comments to 0.77/0.71/0.73 without them. A hybrid detector is then built, combining a comment-independent rule layer for the enumerable patterns with the same classifier gated at a stricter confidence threshold. It restores perfect precision (1.00) in the code-only setting -- matching the rich-context ideal -- but at a real, disclosed recall cost (0.48 vs. the plain classifier's 0.71): the "hard" pattern is by construction unsolvable without the missing comment, and no rule/ML combination recovers information that was never present in the code. Sweeping the confidence threshold (0.5 to 0.9) shows this is a tunable precision/recall trade-off, not a single fixed outcome.
\end{abstract}
